\documentclass[11pt]{article}
\usepackage[margin=0.75in]{geometry}
\usepackage{amsmath,amssymb}
\usepackage{booktabs}
\usepackage{enumitem}
\usepackage{microtype}
\title{XC330 Simulator: Proposed Secondary-Goal Torque-Assist Architecture}
\author{Implementation-derived reference}
\date{}

\begin{document}
\maketitle

This note describes the proposed control architecture for
\texttt{OperatingMode.CURRENT\_POSITION} and
\texttt{Objective.SECONDARY\_GOAL\_ASSIST} are selected. It is a simplified model,
not a numerical reproduction of XC330 firmware. The simulator now exposes this
architecture as its default secondary-goal objective; the older dynamic-current
cap objective remains available for comparison.

\section*{What controls what}
\begin{center}
\begin{tabular}{p{0.25\linewidth} p{0.65\linewidth}}
\toprule
Signal or parameter & Responsibility \\
\midrule
$\hat{\tau}_{ext}$, bias & Motor-side estimate and zeroing; determines the debiased effort $r$. \\
$v_{goal}$ & A filtered, signed intent velocity derived from sustained debiased effort. It controls activation and release. \\
$q_{goal}$ & Persistent secondary goal position, integrated from $v_{goal}$. It is not recomputed as $q+\text{lead}$ each frame. \\
$\ell$ and $q_{ref}$ & Bounded lead ahead of $q_{goal}$; $q_{ref}=q_{goal}+\operatorname{sgn}(v_{goal})\ell$. \\
$\tau_{PID}$ & Desired compliant position-control effort before current limiting; use low $P$, little/no $I$, and $D$ for damping. \\
$I_{ceiling}$, $\tau_{act}$ & Independent current/safety ceiling and actual applied torque after clamping the PID request. \\
\bottomrule
\end{tabular}
\end{center}

\section*{Discrete-time algorithm}
At each simulation step of duration $\Delta t$, the proposed controller executes
the following sequence. The external profile is the simulated plant disturbance;
the observer does not directly read the known disturbance.

\begin{enumerate}[leftmargin=*]
  \item Generate external effort:
  \[
    \tau_{ext,k} = \mu(t_k) + \epsilon_k + \tau_{ext,offset}.
  \]
  \item Estimate external effort from measured acceleration, modeled friction,
  actuator torque, and endstop torque. Low-pass filter the residual to obtain
  $\hat{\tau}_{ext,k}$. During the collection epoch, estimate the bias $b_k$ as
  the running mean. Then compute
  \[
    r_k = \hat{\tau}_{ext,k} - b_k.
  \]
  \item Form a signed intent velocity from the filtered residual:
  \[
    v_{goal,k}^{*} = k_r r_k.
  \]
  Apply a threshold and dwell timer so a short noise excursion does not activate
  assistance. Release when the residual is small or when progress in the
  intended direction remains below a velocity threshold for a release dwell.
  \item Integrate a persistent secondary goal:
  \[
    q_{goal,k+1} = q_{goal,k} + \Delta t\,v_{goal,k}.
  \]
  On release, freeze or capture $q_{goal}$ at the achieved position and decay
  the lead state. This prevents the reference from continuing to run away when
  the user or a physical endstop can no longer produce motion.
  \item Form a bounded target lead and the position reference:
  \[
    q_{ref,k} = q_{goal,k} + \operatorname{sgn}(v_{goal,k})\ell_k,
    \qquad 0 \leq \ell_k \leq \ell_{max}.
  \]
  The soft-endstop position is not used as a clamp on $q_{ref}$. Measured
  velocity is used as a progress signal; if it collapses, lead and assistance
  are released instead.
  \item Compute compliant position-control effort:
  \[
    e_k = q_{ref,k} - q_k,
    \qquad
    \tau_{PID,k} = K_p e_k + K_i\!\int e\,dt
                   + K_d(v_{ref,k}-v_k).
  \]
  \item Apply an independent current/safety ceiling:
  \[
    \tau_{act,k} = \operatorname{clamp}\left(\tau_{PID,k},
      -I_{ceiling}K_t,+I_{ceiling}K_t\right).
  \]
  Here the assist signal changes the target and its lead, while $I_{ceiling}$
  remains an independent compliance/safety parameter.
  \item Advance the plant:
  \[
    \dot v_k = \frac{\tau_{act,k}+\tau_{ext,k}+\tau_{endstop,k}
      -\tau_{friction,k}}{J},
    \qquad
    v_{k+1}=v_k+\Delta t\dot v_k,
    \qquad
    q_{k+1}=q_k+\Delta t v_{k+1}.
  \]
\end{enumerate}

\section*{Interpreting the proposed output stage}
If the view shows $\tau_{PID}=-1.76$ N m and $I_{ceiling}=300$ mA, then the
maximum available torque is
\[
  |\tau_{act}| \le 300\,\mathrm{mA}\times K_t
  \approx 0.345\,\mathrm{N\,m}.
\]
The applied torque is therefore the PID request clipped to about $-0.345$ N m,
not necessarily $-1.76$ N m. The sign is carried by applied current/torque;
the configured Goal Current is a magnitude limit.

\section*{Important design choice}
The proposed design uses the assist estimate to create a secondary goal
velocity and position lead. The position controller remains compliant through
low $P$, minimal $I$, damping $D$, and an independent current ceiling. If
progress stops, assistance releases and the achieved position is captured;
if progress continues, the secondary goal can advance toward the configured
endpoint and remain there.

\section*{Firmware reconciliation (assist-v3)}
The OpenRB implementation in \texttt{assist\_controller.h/.cpp} now implements
this secondary-goal architecture. Defaults are $K_P^{raw}=350$,
$K_I^{raw}=0$, $K_D^{raw}=160$, $I_{ceiling}=300$ mA, $\ell_{max}=8$ degrees,
lead slew 24 degrees/s, detection and stall dwell 0.5 s, progress threshold
1 degree/s and intent conversion 80 degrees/s per N m. The effort-intent
ceiling is 0.25 N m; it is independent of the current ceiling.

The GUI's explicit Enable operation performs grouped selected-ID mode/hold
setup, collects 50 stationary baseline samples at nominal 50 Hz (one second),
and activates automatically. No manual per-joint bias workflow is required.
With requested cap $C$, configured motor limits $L_i$ and total current budget
$B$, installed caps are $\lfloor\min(C,L_i)\min(1,B/\sum_j\min(C,L_j))\rfloor$.
The GUI reports these effective caps; neither safety limits nor the budget are
raised. Insufficient threshold/bias/noise headroom rejects startup. Automatic
activation preserves the host heartbeat deadline rather than renewing it.

All selected integer motor IDs are serviced together on each 20-ms control
pass: one Protocol 2.0 Sync Read of current/velocity/position, one Sync Read of
torque status, then at most one Sync Write containing every changed goal.
There is no round-robin selection or per-ID fallback. Partial replies or status
errors reject the whole frame before any goal write. Motors reply serially on
the bus; a shared frame timestamp does not imply simultaneous physical sampling.
Calibration settings are installed and read back in groups before torque-on;
stopping uses independent zero-current and torque-off Sync Writes followed by
a grouped torque-status read to confirm actual disable. The
firmware observer filters
\[
 r_I^* = (J_d a_d + B_d v_d + C\operatorname{sgn}(v_d))/K_t - I + b_I,
 \qquad \alpha=\frac{\Delta t}{0.20+\Delta t}.
\]
The configurable hardware model terms default to zero, giving a quasi-static
holding-current residual proxy. Simulator plant parameters are illustrative,
not identified motor/linkage dynamics. The simulator's inertia coefficient
multiplies acceleration in degrees/s$^2$; a physical SI inertia requires a
$\pi/180$ conversion for this degree-coordinate equation.
No environmental endstop torque is independently measured or subtracted on
hardware. Environmental resistance is part of its residual load; stalled
progress triggers release. There is no firmware soft-endstop addition.

Stored hard joint limits plus a 2-degree margin bound both the integrated goal
and emitted reference, including inward encoder quantization. Multi-turn wrist
coordinates are preserved. Firmware captures the achieved pose on effort
release as well as stalled progress, and cancels lead immediately on the
release observation. These are intentional refinements: the simulator can
retain lead for one stall-transition sample and decay it against an old goal
on ordinary release. After capture, 0.5 s below half the effort threshold is
required to rearm. The current ceiling remains fixed throughout the session.

Bias collection, per-ID exclusivity, measured-feedback faults, heartbeat lease,
and torque-off retry remain device-side guards. The previous 5-degree bias
window and discrete step/settle gates are removed; a fixed bias still does not
identify gravity throughout ROM. Matching raw PID registers does not make the
simplified simulation a numerical model of the motor's internal controller.
See \texttt{docs/assist\_mode.md} for SDK units, protocol and verification.

\end{document}
